\documentclass[10pt,a4paper]{amsart}
\usepackage[margin=19mm]{geometry}
\usepackage{amsmath,amssymb,mathtools}
\usepackage[scr=rsfs,cal=euler]{mathalfa}
\usepackage{xfrac}
\usepackage[unicode,hidelinks,bookmarks=false]{hyperref}
\newcommand{\ie}{i.e.\ }
\newcommand{\cf}{cf.\ }
\newcommand{\resp}{respectively\ }
\newcommand{\Subsection}{\subsection}
\providecommand{\nobreakdash}{\nobreak\hbox{-}}
\setlength{\emergencystretch}{2em}
\multlinegap0pt
\usepackage{bm}
\usepackage{bbm}
\usepackage{dsfont}
\usepackage{xspace}
\usepackage{cancel}
\usepackage{mathtools}
\usepackage{amsthm}
\makeatletter
\@ifundefined{theorem}{\newtheorem{theorem}{Theorem}}{}
\makeatother
\title[Strong Hamel functions and symmetries]{Strong Hamel functions and symmetries}
\author{Ioan Bucataru, Georgeta Cretu}
\date{Source: arXiv:2402.09791v2 (CC BY 4.0)}
\begin{document}
\maketitle

\noindent\emph{Source and licence.} Strong Hamel functions and symmetries. Version arXiv:2402.09791v2 (CC BY 4.0), distributed under the Creative Commons Attribution 4.0 International licence, as recorded in the arXiv licence metadata for this exact version and shown on the abstract page https://arxiv.org/abs/2402.09791v2. Full rights notice: licenses/arxiv-2402.09791-CC-BY-4.0.txt.\par\smallskip

\noindent\textit{Editorial scope.} Counted unit: the Theorem whose source label is \texttt{thm:SHF} in the article (source0.tex lines 250--257, source label \texttt{thm:SHF}), delivered together with the article's own complete original proof of it (source0.tex lines 258--314, one \texttt{proof} environment of 57 lines). No mathematics is written, repaired or re-derived: statement and proof are the article's own text line for line. Required context retained uncounted: alpha (lines 217--219); ixalpha (lines 223--225). Every definition, display and label of those lines is retained unchanged. Omitted from this package and not redistributed: 1--216, 220--222, 226--249, 315--470. Nothing inside a retained excerpt is omitted or altered: every retained body is delivered exactly as the article's lines are written, so no omission receipt of any kind accompanies this package. Citations: the retained excerpts carry no citation command at all, so no bibliography entry of the article is transcribed and no reference is redistributed. Reference adaptation: none was needed and none was made; every reference inside the delivered unit resolves inside it, so no unresolved reference or citation is produced. Printed slips kept literal: no printed word, symbol, hypothesis or number of the counted statement or of its proof was altered. Layout and class: the article's own class is \texttt{amsart} with packages including babel, amsthm, gensymb, hyperref, multicol, extarrows, mathtools, commath, amssymb, amscd; the wrapper is the lane's pinned \texttt{amsart} editorial wrapper at 10pt on A4, which restates the theorem-like environments the retained text needs and adds the article's own macro definitions that the retained text needs (none), with extra wrapper packages (bm, bbm, dsfont, xspace, cancel, mathtools). The delivered numbering is the unit's own. Assets: no retained span contains a figure environment, a graphics-inclusion command, a PDF-page inclusion, a table or a TikZ picture; no image file is copied into this package, the package declares no asset dependency and no image file is redistributed with it. Licence: version arXiv:2402.09791v2 of this article is distributed under the Creative Commons Attribution 4.0 International licence, as recorded in the arXiv licence metadata for this exact version and shown on the abstract page https://arxiv.org/abs/2402.09791v2. A full rights notice is delivered at licenses/arxiv-2402.09791-CC-BY-4.0.txt. Characters: every retained excerpt is pure ASCII, so the delivered engine, XeLaTeX with its default Latin Modern text font, reports no missing character.
\begin{equation}\label{alpha}
\alpha = i_{[J,G]}\mathcal{L}_Gd_Jf' = \nabla\left(\dfrac{\partial f'}{\partial y^i}\right)dx^i-\dfrac{\partial f'}{\partial y^i}\delta y^i = \left(G\left(\dfrac{\partial f'}{\partial y^i}\right) - 2N^j_i\dfrac{\partial f'}{\partial y^j}\right)dx^i-\dfrac{\partial f'}{\partial y^i} d y^i.
\end{equation} 

\begin{eqnarray}
i_Xdd_JL=\alpha. \label{ixalpha}
\end{eqnarray}

\begin{theorem}\label{thm:SHF}
We consider $G$ the geodesic spray of a Finsler metric $F$. Then, the following statements are equivalent:
\begin{itemize}
\item[i)] There exists a strong Hamel function for the geodesic spray $G$.
\item[ii)] There exists a strong dual symmetry for $G$.
\item[iii)] There exists a strong dynamical symmetry for $G$.
\end{itemize}
\end{theorem}

\begin{proof}
We will prove the equivalence between $i)$ and $ii)$. If $f$ is a strong Hamel function of the geodesic spray $G$, then there exists a $0$-homogeneous potential function $f'$ such that $f=G(f')=\nabla f'$ and $\delta_Gf=0$.

Using the commutation rule $d_J\nabla=\nabla d_J+d_h+2i_R$, we obtain the following equivalent expression of the $1$-form $\alpha$ introduced in \eqref{alpha}:
\begin{eqnarray}
\alpha=\nabla d_Jf'-d_vf'=d_J\nabla f'-d_hf'-d_vf'=d_Jf-df'. \label{alphaff'}
\end{eqnarray} 
We can now verify the geodesic invariance of the $1$-form $\alpha$:
\begin{eqnarray*}
\mathcal{L}_{G}\alpha=\mathcal{L}_{G}d_Jf-\mathcal{L}_{G}df'=\mathcal{L}_{G}d_Jf-dG(f')=\mathcal{L}_{G}d_Jf-df=\delta_Gf=0.
\end{eqnarray*}
From the last expression \eqref{alphaff'} of the $1$-form $\alpha$, we obtain $i_J\alpha=-d_Jf'$, which means that $\alpha$ is a strong dual symmetry.

For the converse, we consider $\alpha$ a $0$-homogeneous $1$-form on $T_0M$, which can be written as
$ \alpha=\alpha_i dx^i+\beta_i dy^i$. We now assume that $i_J\alpha = \beta_i dx^i$ is $d_J$-exact, and hence there exists a $0$-homogeneous function $f'(x,y)$ such that $i_J\alpha=-d_Jf'$.
Therefore, the $1$-form $\alpha$ is given by
\begin{eqnarray} \label{alphaf'}
\alpha=\alpha_i dx^i - \frac{\partial f'}{\partial y^i}dy^i.
\end{eqnarray}
 The Lie derivative of the $1$-form \eqref{alphaf'} along the geodesic spray $G$ is given by:
\begin{eqnarray*}
\mathcal{L}_G\alpha = G(\alpha_i)dx^i + \alpha_i dy^i - G\left(\frac{\partial f'}{\partial y^i}\right) dy^i + \frac{\partial f'}{\partial y^i} 2\left(\frac{\partial G^i}{\partial x^j}dx^j + \frac{\partial G^i}{\partial y^j} dy^j\right).  
\end{eqnarray*}
If the $1$-form $\alpha$ is geodesically invariant, then collecting the vertical components, we obtain 
\begin{eqnarray*}
\alpha_i - G\left(\frac{\partial f'}{\partial y^i}\right) + 2 \frac{\partial f'}{\partial y^j} \frac{\partial G^j}{\partial y^i}=0. 
\end{eqnarray*}
We now use the commutator: 
\begin{eqnarray*}
\left[G, \frac{\partial}{\partial y^i}\right] = \frac{\partial}{\partial x^i} - 2N^k_i\frac{\partial}{\partial y^k},
\end{eqnarray*}
to obtain the components $\alpha_i$:
\begin{eqnarray*}
\alpha_i = G\left(\frac{\partial f'}{\partial y^i}\right) - 2 \frac{\partial f'}{\partial y^j} \frac{\partial G^j}{\partial y^i} = \frac{\partial}{\partial y^i}\left(Gf'\right) + \frac{\partial f'}{\partial x^i}.
\end{eqnarray*}
The function $f:=G(f')$ is $1$-homogeneous and the $1$-form $\alpha$ is given by:
\begin{eqnarray*}
\alpha=\left(\frac{\partial f}{\partial y^i} -  \frac{\partial f'}{\partial x^i}\right) dx^i - \frac{\partial f'}{\partial y^i}dy^i = d_Jf-df',
\end{eqnarray*} 
 which coincides with formula \eqref{alphaff'}. Therefore, we have $0=\mathcal{L}_G\alpha=\delta_Gf$ and hence $f=G(f')$ is a strong Hamel function. 
 
We note that a strong Hamel function and the corresponding strong dual symmetry have the same potential function. 

Now, we prove the equivalence between $ii)$ and $iii)$. For a $1$-form $\alpha\in \bigwedge^1(T_0M)$ we consider $X$ its symplectic dual vector field, given by \eqref{ixalpha}. Therefore, we have
\begin{eqnarray*} 
\mathcal{L}_G\alpha = \mathcal{L}_Gi_Xdd_JL=i_X\mathcal{L}_Gdd_JL +i_{[G,X]}dd_JL = i_{[G,X]}dd_JL. 
\end{eqnarray*}
Since $dd_JL$ is a symplectic form, we have that $\mathcal{L}_G\alpha=0$ if and only if $[G,X]=0$, which means that $\alpha$ is a dual symmetry for the geodesic spray $G$ if and only if $X$ is a dynamical symmetry of $G$. We now show that the dual symmetry $\alpha$ is strong if and only if its dual vector field is a strong dynamical symmetry, with the same potential function. 

For a vector field, $X=X^i{\partial}/{\partial x^i} + Y^i{\partial}/{\partial y^i}\in \mathfrak{X}(T_0M)$, homogeneous of order $-1$, $JX=X^i{\partial}/{\partial y^i}$ is a vertical gradient if there exists a $0$-homogeneous function $f'$ on $T_0M$ such that $g_{ij}X^j=\frac{\partial f'}{\partial y^i}$. 

We start form the duality relation between the $1$-form $\alpha$ and the corresponding vector field $X$, \eqref{ixalpha}, we apply the algebraic derivation $i_J$, and use the fact that $i_Jdd_JL=0$. Hence, we have:
\begin{eqnarray*}
i_J\alpha=i_Ji_Xdd_JL=i_Xi_Jdd_JL - i_{JX}dd_JL=-i_{JX}dd_JL= - i_{X^k\frac{\partial}{\partial y^k}}g_{ij}\delta y^i\wedge dx^j = - g_{ij}X^idx^j.
\end{eqnarray*} 
Therefore $i_J\alpha$ is $d_J$-exact with potential function $-f'$ if and only if $JX$ is a gradient with potential function $f'$. This completes the proof of the equivalence between ii) and iii).   
\end{proof}

\end{document}
